\chapter{Effective Field Theory}

\begin{remarkbox}
Effective field theory is the modern language of QFT at finite resolution. It
accepts that a theory may be valid only below some scale and organizes all
allowed interactions as an expansion in powers of energy over that scale.
\end{remarkbox}

\section{The Effective Field Theory Viewpoint}

The renormalization group teaches that physics changes with scale. Effective
field theory, or EFT, turns this lesson into a practical method. Instead of
requiring a theory to be fundamental at all distances, one asks what degrees of
freedom and symmetries are relevant at the energies being studied.

An EFT is valid below a cutoff scale \(\Lambda\). It contains the fields that are
light compared with \(\Lambda\), and it includes every interaction compatible
with the symmetries:
\[
    \mathcal{L}_{\mathrm{EFT}}
    = \mathcal{L}_{\mathrm{ren}}
      +\sum_i \frac{c_i}{\Lambda^{\Delta_i-4}}\mathcal{O}_i.
\]
Here \(\mathcal{O}_i\) are local operators of dimension \(\Delta_i>4\), and
\(c_i\) are dimensionless coefficients.

The expansion parameter is usually
\[
    \frac{E}{\Lambda},
\]
where \(E\) is the characteristic energy of the process.

\section{Separation of Scales}

EFT works when there is a hierarchy of scales. Suppose a theory contains a heavy
particle of mass \(M\) and we study processes with
\[
    E\ll M.
\]
At such energies, the heavy particle cannot be produced on shell. Its effects
are still present, but they appear indirectly through local interactions among
the light fields.

The heavy propagator illustrates the idea. For momentum transfer \(p^2\ll M^2\),
\[
    \frac{i}{p^2-M^2}
    = -\frac{i}{M^2}\left(1+\frac{p^2}{M^2}+\frac{p^4}{M^4}+\cdots\right).
\]
The nonlocal exchange of a heavy particle becomes a series of local operators
suppressed by powers of \(M\).

This is the basic mechanism behind effective interactions.

\section{Integrating Out Heavy Fields}

To integrate out a heavy field means to remove it as an explicit degree of
freedom while keeping its low-energy effects. In the path integral, if \(H\) is a
heavy field and \(\ell\) denotes light fields, one writes schematically
\[
    e^{iS_{\mathrm{eff}}[\ell]}
    = \int\mathcal{D}H\,e^{iS[\ell,H]}.
\]
The result is an effective action for the light fields only.

At tree level, integrating out a heavy field often amounts to solving its
classical equation of motion and substituting back. At loop level, it also
produces determinants and quantum corrections.

The effective action is generally nonlocal if written exactly. At low energies,
it is expanded in local operators with increasing numbers of derivatives or
fields.

\section{Matching}

Matching determines the coefficients of the EFT operators. One computes the same
low-energy quantity in the full theory and in the EFT, then chooses the EFT
coefficients so that the two agree at the matching scale.

Schematically,
\[
    \mathcal{A}_{\mathrm{full}}(E\ll M)
    = \mathcal{A}_{\mathrm{EFT}}(E\ll M).
\]
This equality fixes the Wilson coefficients \(c_i\).

Matching can be done at tree level or loop level. Tree-level matching captures
classical heavy-particle exchange. Loop-level matching captures virtual heavy
particle effects.

The matching scale is usually chosen near the heavy mass:
\[
    \mu\sim M.
\]
Then large logarithms are avoided at the matching step. Renormalization group
evolution can then run the EFT coefficients down to lower scales.

\section{Power Counting}

Power counting tells us which operators matter at a given accuracy. If an
operator has dimension \(\Delta\), its contribution scales roughly as
\[
    \left(\frac{E}{\Lambda}\right)^{\Delta-4}.
\]
Operators of higher dimension are more suppressed at low energies.

This means that EFT can be predictive even though it contains infinitely many
operators. To compute to a fixed order in \(E/\Lambda\), only finitely many
operators are needed.

For example, if one works up to order \(1/\Lambda^2\), only dimension-six
operators are needed beyond the renormalizable terms. Dimension-eight operators
enter at order \(1/\Lambda^4\).

\section{Symmetries and Operator Bases}

The EFT Lagrangian must include all operators allowed by the symmetries. It must
exclude operators forbidden by symmetries. Symmetries therefore organize the
operator basis.

Operators related by integration by parts, equations of motion, or algebraic
identities are not independent. An operator basis is a minimal independent set of
operators after these redundancies are removed.

The construction follows a clear logic:
\begin{enumerate}
    \item choose the light fields;
    \item choose the symmetries;
    \item list local operators ordered by dimension;
    \item remove redundant operators;
    \item assign Wilson coefficients;
    \item match coefficients to data or to a UV theory.
\end{enumerate}
The result is the most general low-energy theory compatible with the chosen
assumptions.

\section{Wilson Coefficients}

The coefficients \(c_i\) multiplying EFT operators are called Wilson
coefficients. They summarize short-distance physics not resolved explicitly in
the EFT.

A Wilson coefficient may be computed from a known high-energy theory by
matching. If the high-energy theory is unknown, the coefficients are treated as
parameters to be measured or bounded experimentally.

The separation is
\[
    \text{long-distance physics} \longrightarrow \text{matrix elements of EFT operators},
\]
\[
    \text{short-distance physics} \longrightarrow \text{Wilson coefficients}.
\]
Renormalization group evolution moves information between scales by changing the
coefficients as \(\mu\) changes.

\section{Decoupling}

The decoupling principle says that heavy particles have suppressed effects at low
energies. If a heavy particle has mass \(M\), its virtual effects often appear as
powers of \(1/M\).

However, decoupling must be interpreted carefully. Heavy particles can also
modify renormalized parameters of lower-dimensional operators. Such effects are
not necessarily suppressed by powers of \(1/M\); they may appear as threshold
corrections to masses, couplings, or kinetic terms.

After these parameters are matched, genuinely new low-energy effects from heavy
physics are organized by higher-dimensional operators.

\section{Naturalness}

Naturalness asks whether small parameters remain small under quantum corrections.
Suppose a scalar mass receives a correction of order
\[
    \delta m^2\sim \frac{\lambda}{16\pi^2}\Lambda^2.
\]
If the physical mass is much smaller than \(\Lambda\), a cancellation between the
bare mass and the correction may be required. This is called fine-tuning.

Fermion masses and gauge boson masses can be protected by symmetries. Scalar
masses are less protected in generic theories. This is the origin of naturalness
questions surrounding scalar fields, including the Higgs boson.

Naturalness is not a theorem. It is a diagnostic about sensitivity to short-
distance physics.

\section{The Fermi Theory of Weak Interactions}

A classic EFT example is Fermi's four-fermion theory. At energies much below the
\(W\)-boson mass, weak interactions can be described by a local operator:
\[
    \mathcal{L}_{\mathrm{Fermi}}
    \sim -\frac{G_F}{\sqrt{2}}(\bar\psi\gamma^\mu(1-\gamma^5)\psi)
    (\bar\psi\gamma_\mu(1-\gamma^5)\psi).
\]
The coefficient scales as
\[
    G_F\sim \frac{1}{M_W^2}.
\]
The full electroweak theory contains \(W\)-boson exchange. At low energies,
that exchange becomes a local four-fermion interaction.

This example shows how an EFT can be extremely successful while not being
fundamental at arbitrarily high energies.

\section{Euler--Heisenberg as a Photon EFT}

Another example is the low-energy effective theory of photons after integrating
out the electron. At energies much smaller than the electron mass, virtual
electron loops generate photon self-interactions:
\[
    \mathcal{L}_{\mathrm{EH}}
    = -\frac{1}{4}F_{\mu\nu}F^{\mu\nu}
      + c_1\frac{(F_{\mu\nu}F^{\mu\nu})^2}{m_e^4}
      + c_2\frac{(F_{\mu\nu}\widetilde F^{\mu\nu})^2}{m_e^4}
      +\cdots.
\]
These terms describe light-by-light scattering at energies below the electron
mass. They are strongly suppressed by powers of \(1/m_e^4\).

The EFT makes clear why photon self-interactions are tiny at low energies even
though they are allowed quantum mechanically.

\section{EFT and Renormalization}

EFTs are renormalized like any other quantum field theory. Loop diagrams
involving higher-dimensional operators generate divergences, which are absorbed
into coefficients of operators allowed by the symmetries.

Because infinitely many operators are present, the theory is not renormalizable
in the old finite-counterterm sense. But it is renormalizable order by order in
\(E/\Lambda\). At any fixed order, only finitely many counterterms are needed.

This is the modern meaning of predictivity in EFT: choose an accuracy, include
all operators up to the required dimension, renormalize them, and estimate the
size of neglected terms.

\section{EFT as the Default View of QFT}

Modern QFT often treats every theory as effective unless proven otherwise. The
Standard Model itself may be an EFT below a much higher scale. Gravity is also
successfully described as an EFT at energies far below the Planck scale.

This viewpoint changes the meaning of nonrenormalizable interactions. They are
not automatically forbidden. They are suppressed by a high scale and become
important when experiments reach enough precision or energy.

The EFT attitude is therefore practical and humble:
\[
    \text{write what is allowed, organize by scale, measure what is unknown}.
\]

\begin{summarybox}
Effective field theory describes physics below a cutoff scale \(\Lambda\) using
light fields, symmetries, and a systematic expansion in \(E/\Lambda\). Heavy
fields are integrated out, leaving local operators with Wilson coefficients.
Matching fixes those coefficients, power counting determines which operators are
needed, and renormalization group evolution moves the theory between scales. EFT
is predictive order by order even when infinitely many operators are allowed.
\end{summarybox}

\section{Chapter Checklist}

After reading this chapter, one should be able to:
\begin{itemize}
    \item explain the EFT viewpoint;
    \item describe separation of scales;
    \item explain what it means to integrate out a heavy field;
    \item define matching;
    \item use power counting to rank operators;
    \item explain the role of symmetries in building an operator basis;
    \item define Wilson coefficients;
    \item explain decoupling and threshold effects;
    \item describe naturalness as sensitivity to high scales;
    \item explain the Fermi theory as an EFT;
    \item state why EFTs are predictive order by order.
\end{itemize}
